

\documentclass[12pt,a4paper,oneside]{book}

\usepackage{ubthesis}

\usepackage[abbreviations,nomain,nopostdot,nogroupskip]{glossaries-extra}
\setabbreviationstyle[abbreviation]{long-short}


\newabbreviation{abac}{ABAC}{Attribute-Based Access Control}
\newabbreviation{ack}{ACK}{Acknowledgement}
\newabbreviation{adc}{ADC}{Analogue-to-Digital Converter}
\newabbreviation{api}{API}{Application Programming Interface}
\newabbreviation{arm64}{ARM64}{64-bit Advanced RISC Machine architecture}

\newabbreviation{batman}{BATMAN-adv}{Better Approach To Mobile Ad-hoc Networking, advanced}
\newabbreviation{ble}{BLE}{Bluetooth Low Energy}
\newabbreviation{bft}{BFT}{Byzantine Fault Tolerance}

\newabbreviation{ca}{CA}{Certificate Authority}
\newabbreviation{cap}{CAP}{Consistency, Availability, Partition tolerance}
\newabbreviation{cpu}{CPU}{Central Processing Unit}
\newabbreviation{crt}{CRT}{Chinese Remainder Theorem}

\newabbreviation{did}{DID}{Decentralised Identifier}
\newabbreviation{dlt}{DLT}{Distributed Ledger Technology}
\newabbreviation{dp}{DP}{Differential Privacy}

\newabbreviation{edp}{EDP}{Energy-Delay Product}
\newabbreviation{esp32}{ESP32}{Espressif 32-bit microcontroller with integrated Wi-Fi and Bluetooth}

\newabbreviation{gdpr}{GDPR}{General Data Protection Regulation}
\newabbreviation{gw}{GW}{Gateway}

\newabbreviation{hlf}{HLF}{Hyperledger Fabric}
\newabbreviation{hrbac}{HRBAC}{Hierarchical Role-Based Access Control}

\newabbreviation{iot}{IoT}{Internet of Things}
\newabbreviation{ipfs}{IPFS}{InterPlanetary File System}

\newabbreviation{json}{JSON}{JavaScript Object Notation}

\newabbreviation{lora}{LoRa}{Long Range (sub-GHz spread-spectrum radio modulation)}
\newabbreviation{lorawan}{LoRaWAN}{Long Range Wide Area Network}
\newabbreviation{lzw}{LZW}{Lempel-Ziv-Welch compression}

\newabbreviation{mac}{MAC}{Medium Access Control}
\newabbreviation{mcu}{MCU}{Microcontroller Unit}
\newabbreviation{msp}{MSP}{Membership Service Provider}

\newabbreviation{pbft}{PBFT}{Practical Byzantine Fault Tolerance}
\newabbreviation{pki}{PKI}{Public Key Infrastructure}

\newabbreviation{ram}{RAM}{Random Access Memory}
\newabbreviation{rbac}{RBAC}{Role-Based Access Control}
\newabbreviation{rf}{RF}{Radio Frequency}
\newabbreviation{roi}{ROI}{Return on Investment}
\newabbreviation{rpi}{RPi}{Raspberry Pi}
\newabbreviation{rtc}{RTC}{Real-Time Clock}

\newabbreviation{sf}{SF}{Spreading Factor}
\newabbreviation{sha}{SHA}{Secure Hash Algorithm}
\newabbreviation{slo}{SLO}{Service Level Objective}
\newabbreviation{sss}{SSS}{Shamir Secret Sharing}

\newabbreviation{tee}{TEE}{Trusted Execution Environment}
\newabbreviation{tps}{TPS}{Transactions Per Second}
\newabbreviation{tpm}{TPM}{Trusted Platform Module}

\newabbreviation{zkp}{ZKP}{Zero-Knowledge Proof}
 
\usepackage[
  backend=biber,
  style=numeric-comp,
  sorting=none,
  maxbibnames=6,
  minbibnames=3,
  giveninits=true,
  doi=true, isbn=false, url=false
]{biblatex}
\addbibresource{references.bib}

\begin{document}

\frontmatter
\pagestyle{fancy}



\thispagestyle{empty}
\begin{singlespace}
\fontsize{14}{16.1}\selectfont

\vspace*{-0.9\baselineskip}

\begin{center}
  \textbf{UNIVERSITY OF BUEA}
\end{center}

\vspace{3.72\baselineskip}

\noindent
\begin{minipage}[t]{0.56\textwidth}
  \raggedright
  \textbf{FACULTY OF ENGINEERING\\AND TECHNOLOGY}
\end{minipage}\begin{minipage}[t]{0.44\textwidth}
  \raggedright
  \textbf{DEPARTMENT OF\\COMPUTER ENGINEERING}
\end{minipage}

\vspace{4.59\baselineskip}

\begin{center}
\textbf{A DISTRIBUTED BLOCKCHAIN-BASED APPROACH TO\\
          PROTECT IDENTITY, INTEGRITY AND PRIVACY IN\\
          AGRICULTURAL DATA PROCESSING}

  \vspace{4.34\baselineskip}

  {\fontsize{12}{14}\selectfont By}

  \vspace{3.56\baselineskip}

  \textbf{FORCHA GLEN BELOA}\\
  (FE22P049)\\
  M.Eng. in Telecommunications and Networks

  \vspace{2.52\baselineskip}

  A Thesis submitted to the Faculty of Engineering and Technology of the\\
  University of Buea in partial fulfilment of the requirements\\
  for the degree of Doctor of Philosophy (PhD)\\
  in Computer Engineering.
\end{center}

\vspace{2.52\baselineskip}

\noindent\textbf{Supervisors}

\vspace{1.0\baselineskip}

\noindent\textbf{Michael Ekonde SONE, PhD}\\
Associate Professor

\vspace{0.8\baselineskip}

\noindent\textbf{Elie FUTE TAGNE, PhD}\\
Associate Professor

\vfill

\begin{center}
  \textbf{\monthyear}
\end{center}

\vspace{1.5\baselineskip}

\end{singlespace}
\clearpage
 
\clearpage\thispagestyle{empty}\mbox{}   \clearpage


\thispagestyle{empty}
\begin{singlespace}
\fontsize{14}{16.1}\selectfont

\vspace*{-0.9\baselineskip}

\begin{center}
  \textbf{UNIVERSITY OF BUEA}
\end{center}

\vspace{3.72\baselineskip}

\noindent
\begin{minipage}[t]{0.56\textwidth}
  \raggedright
  \textbf{FACULTY OF ENGINEERING\\AND TECHNOLOGY}
\end{minipage}\begin{minipage}[t]{0.44\textwidth}
  \raggedright
  \textbf{DEPARTMENT OF\\COMPUTER ENGINEERING}
\end{minipage}

\vspace{4.59\baselineskip}

\begin{center}
\textbf{A DISTRIBUTED BLOCKCHAIN-BASED APPROACH TO\\
          PROTECT IDENTITY, INTEGRITY AND PRIVACY IN\\
          AGRICULTURAL DATA PROCESSING}

  \vspace{4.34\baselineskip}

  {\fontsize{12}{14}\selectfont By}

  \vspace{3.56\baselineskip}

  \textbf{FORCHA GLEN BELOA}\\
  (FE22P049)\\
  M.Eng. in Telecommunications and Networks

  \vspace{2.52\baselineskip}

  A Thesis submitted to the Faculty of Engineering and Technology of the\\
  University of Buea in partial fulfilment of the requirements\\
  for the degree of Doctor of Philosophy (PhD)\\
  in Computer Engineering.
\end{center}

\vspace{2.52\baselineskip}

\noindent\textbf{Supervisors}

\vspace{1.0\baselineskip}

\noindent\textbf{Michael Ekonde SONE, PhD}\\
Associate Professor

\vspace{0.8\baselineskip}

\noindent\textbf{Elie FUTE TAGNE, PhD}\\
Associate Professor

\vfill

\begin{center}
  \textbf{\monthyear}
\end{center}

\vspace{1.5\baselineskip}

\end{singlespace}
\clearpage
 

\setcounter{page}{2}                     

\ubfrontheading{Dedication}
\vspace{2\baselineskip}
\begin{center}
To \dots
\end{center}
\clearpage
 \ubfrontheading{Certification}
\begin{singlespace}
\vspace{\baselineskip}
\begin{center}{\bfseries UNIVERSITY OF BUEA}\\
{\bfseries FACULTY OF ENGINEERING AND TECHNOLOGY}\end{center}
\vspace{2\baselineskip}

This is to certify that the work described in the thesis entitled
\textbf{``A Distributed Blockchain-Based Approach to Protect Identity,
Integrity and Privacy in Agricultural Data Processing''}

\vspace{\baselineskip}
By \hrulefill

\vspace{\baselineskip}
Was carried out in the Department of Computer Engineering

\vspace{\baselineskip}
Under the supervision of

\vspace{2\baselineskip}
\hrulefill\\
(Name and Signature of Supervisor)

\vspace{2\baselineskip}
Date \hrulefill
\end{singlespace}
\clearpage
 \ubfrontheading{Acknowledgements}

\clearpage
 \ubfrontheading{Abstract}

\clearpage
 
\ubfrontheading{Table of Contents}
\begin{singlespace}
\renewcommand{\contentsname}{}
\vspace{-3\baselineskip}
\tableofcontents
\end{singlespace}
\clearpage

\ubfrontheading{List of Figures}
\begin{singlespace}
\renewcommand{\listfigurename}{}
\vspace{-3\baselineskip}
\listoffigures
\end{singlespace}
\clearpage

\ubfrontheading{List of Tables}
\begin{singlespace}
\renewcommand{\listtablename}{}
\vspace{-3\baselineskip}
\listoftables
\end{singlespace}
\clearpage

\ubfrontheading{List of Abbreviations}
\begin{singlespace}
\printunsrtglossary[type=abbreviations,title={},style=long]
\end{singlespace}
\clearpage

\mainmatter

\clearpage{}

\chapter{Introduction}
\label{ch:introduction}

\section{Background to the Study}
\label{sec:background}

Agriculture is undergoing a transition from experience-driven practice to
data-driven decision making. Networks of low-cost sensors now measure soil
moisture, canopy temperature, humidity and nutrient status continuously, and
the resulting streams feed irrigation controllers, yield models and
certification schemes~\cite{rathi2025smart, brooks2019precisionWater}. In
water-stressed regions this shift is not a convenience but a necessity.
Irrigation accounts for the dominant share of freshwater withdrawal, and
climate variability has compressed the margin for error in scheduling
decisions~\cite{miller2020climateImpact, williams2022savings}.

A measurement taken in a field does not stop at the field. It becomes a record,
and that record is consumed by parties with interests of their own, as
Figure~\ref{fig:lifecycle} sets out. The value of the record to those parties
is exactly what creates the incentive to falsify it, and the further a record
travels from the instrument that produced it, the weaker the evidence that it
is genuine.

\begin{figure}[htbp]
  \centering
  \begin{tikzpicture}[
      box/.style={draw, rounded corners=2pt, align=center, font=\small,
                  minimum height=9mm, inner sep=4pt},
      use/.style={box, fill=black!4, text width=26mm},
      flow/.style={-{Stealth[length=2mm]}, thick},
      note/.style={font=\footnotesize\itshape, align=center, text width=90mm}
    ]
    \node[box, text width=23mm] (field)  at (0,0)      {Field\\measurement};
    \node[box, text width=23mm] (record) at (3.7,0)    {Agricultural\\record};

    \node[use] (cert) at (8.0, 1.5)  {Certification\\(organic, fair trade)};
    \node[use] (ins)  at (8.0, 0)    {Index-based\\crop insurance};
    \node[use] (sub)  at (8.0,-1.5)  {Subsidy\\allocation};

    \draw[flow] (field) -- (record);
    \draw[flow] (record.east) -- (cert.west);
    \draw[flow] (record.east) -- (ins.west);
    \draw[flow] (record.east) -- (sub.west);

\node[note] (inc) at (4.0,-2.9)
      {Each downstream use attaches value to the record,\\
       and therefore creates an incentive to falsify it.};
  \end{tikzpicture}
  \caption{Lifecycle of an agricultural measurement. A reading becomes a record
           that is consumed by parties outside the farm, and the value those
           parties attach to it is what makes falsification worthwhile.}
  \label{fig:lifecycle}
\end{figure}

The economics of the sensing layer explain why the problem is hard rather than
merely tedious. A deployment that is worth installing on a smallholding must
cost a small fraction of the crop value it protects, which forces the node bill
of materials down to a level at which neither a hardware security module nor a
generous energy budget is affordable. The same economics rule out mains power
and wired backhaul across a field, so the design space collapses onto
battery-powered microcontrollers speaking a long-range, low-rate radio. Every
protection mechanism proposed in this thesis has to survive inside that
envelope, and a mechanism that works only on a mains-powered gateway has solved
the easier half of the problem.

The sensing layer of such systems is built from severely constrained devices.
An \gls{esp32} node operating on battery and harvested solar power must
transmit over kilometre-scale distances using \gls{lora}, a modulation scheme
that trades data rate for range and therefore charges a high airtime cost per
packet~\cite{raza2018lorawan, adelantado2017understanding}. Airtime is the
binding constraint. It determines energy consumption at the node, and, through
duty-cycle regulation and collision probability, it determines how many nodes a
single gateway can serve~\cite{mohammed2025loraEnergy}. Conventional responses
compress the payload, but general-purpose compressors such as Huffman and
\gls{lzw} carry dictionary and buffer overheads that a microcontroller with a
few kilobytes of free \gls{ram} can ill
afford~\cite{marcelloni2009energy, chen2025compression}.

Two further constraints sharpen this. Sub-GHz \gls{iot} bands are governed by a
duty-cycle limit, so a node that has just transmitted must stay silent for a
period proportional to the airtime it consumed. Airtime is therefore not only
an energy cost but a rate limit: it caps how often a node may report, and no
increase in battery capacity relaxes it. Second, because \gls{lora} uses
unslotted access, packets collide whenever two nodes transmit on the same
channel at the same time, and the collision probability grows with the square
of offered traffic in the regime that matters. Halving per-packet occupancy
therefore buys more than a proportional reduction in collisions, which is why
the mechanism studied here is worth the arithmetic it costs.

The \gls{crt} offers a different route. Decomposing a reading into residues
modulo pairwise coprime moduli yields several short packets in place of one
long one~\cite{garner1959residue, ding1996remainder}. Because a multi-channel
concentrator such as the SX1302 demodulates channels concurrently, the short
residue packets can be received simultaneously, and the original value is
recovered in $O(k)$ operations for $k$ moduli. The practical effect is a
reduction in per-packet occupancy rather than in total transmitted volume,
which is precisely the quantity that governs gateway capacity. The detailed
encoding and reconstruction path is presented in Chapter Three.

Residue arithmetic is not new to distributed systems. The same decomposition has
been applied to the throughput problem in permissionless blockchains, where
splitting work across residue classes permits parallel processing of operations
that would otherwise serialise. That line of work establishes the arithmetic and
the parallelisation argument; what it does not address is the constrained radio
link, where the payoff is airtime rather than processor cycles, nor the identity
and confidentiality questions that arise when the residues cross a shared
medium. This thesis takes the arithmetic as established and asks what it is
worth on the edge of the network rather than at its centre.



Data that has been collected cheaply must then be trusted. Agricultural records
increasingly carry consequences beyond the farm gate. They substantiate organic
and fair-trade certification, they underwrite index-based crop insurance, and
they inform subsidy allocation. Each of these uses creates an incentive to
falsify. A centralised database concentrates that risk in a single
administrator, and an unprotected sensor network invites the injection of
fabricated readings by any party with a compatible radio. \gls{dlt} addresses
the tamper-evidence requirement by binding records into a hash-linked chain
replicated across mutually distrusting parties, and permissioned platforms such
as \gls{hlf} add an identity layer through a \gls{msp} and \gls{pki}, making it
feasible to say not merely that a record is unaltered but who produced
it~\cite{tar2023fabric, jayaram2023hyperledger, dube2024fabricIrrigation}.

Running that machinery at the edge is the difficulty. Consensus, endorsement and
world-state maintenance were designed for datacentre hardware, whereas an
agricultural deployment must place gateways in the field, on \gls{arm64}
single-board computers, behind intermittent connectivity. Whether a permissioned
ledger can meet irrigation-control timing requirements on such hardware has
until recently been asserted rather than measured.

Finally, protection is not exhausted by integrity. A soil-moisture trace is also
a disclosure. It reveals plot productivity, irrigation strategy and, in
aggregate, the commercial position of the farmer. Where records are replicated
across a consortium that includes buyers, insurers and cooperatives, the same
replication that secures integrity enlarges the audience for private
information. Identity, integrity and privacy therefore have to be treated as one
design problem rather than three independent ones, which is the position this
thesis takes.

\section{Statement of the Problem}
\label{sec:problem}

Three properties that ought to hold simultaneously are, in current systems,
traded against one another.

\textbf{Identity is weak at the sensing layer.} Constrained nodes seldom carry
credentials strong enough to authenticate individual readings. Where
authentication exists it typically terminates at the gateway, so the ledger
attests to what the gateway submitted rather than to what the sensor measured.
The provenance chain has a gap exactly where falsification is cheapest, at the
final hop between an individual node and the record that claims to represent
it.

\textbf{Integrity guarantees do not survive the move to the edge.} Permissioned
ledgers provide tamper evidence, but their published performance
characterisations assume server-class hosts. On \gls{arm64} gateways with
constrained memory and intermittent links, it is not established that
endorsement and ordering can complete inside the control loop of an irrigation
decision, nor how the ledger behaves during network partition and recovery.

\textbf{Privacy is sacrificed to obtain integrity.} Replicating raw agronomic
readings to every consortium member is the simplest way to make them verifiable
and the worst way to keep them confidential. Existing deployments largely accept
this trade, or push confidential data off chain in a way that reintroduces the
trusted intermediary the ledger was meant to remove.

The three shortfalls, shown together in Figure~\ref{fig:three-properties}, are
not independent faults to be repaired one at a time. Each follows from the same
scarcity, and a remedy applied to one is paid for out of the budget available
to the others.

\begin{figure}[htbp]
  \centering
  \begin{tikzpicture}[
      prop/.style={draw, thick, rounded corners=2pt, align=center,
                   font=\small\bfseries, minimum height=8mm, inner sep=4pt,
                   fill=white},
      gap/.style={font=\footnotesize\itshape, align=center, text width=34mm},
      link/.style={{Stealth[length=2mm]}-{Stealth[length=2mm]}, gray, thin}
    ]
    \node[prop] (id)  at (0,2.5)   {Identity};
    \node[prop] (int) at (-3.2,0)  {Integrity};
    \node[prop] (pri) at (3.2,0)   {Privacy};

    \draw[link] (id)  -- (int);
    \draw[link] (int) -- (pri);
    \draw[link] (pri) -- (id);

\node[gap] (gid)  at (0,3.5)     {authentication stops\\at the gateway};
    \node[gap] (gint) at (-3.2,-1.1) {ledger cost unproven\\on edge hardware};
    \node[gap] (gpri) at (3.2,-1.1)  {confidentiality traded\\away for verifiability};

    \begin{scope}[on background layer]
      \node[draw, dashed, thick, rounded corners=4pt, fill=black!3,
            fit=(id)(int)(pri)(gid)(gint)(gpri), inner sep=5mm] (budget) {};
    \end{scope}
    \node[font=\small\itshape, below=1.5mm of budget.south]
      {bounded by one shared energy, memory and channel budget};
  \end{tikzpicture}
  \caption{The three protection properties and the shortfall in each. Because
           all three draw on a single constrained budget, they form one design
           problem rather than three independent ones.}
  \label{fig:three-properties}
\end{figure}

Underlying all three is a resource problem. The transmission budget that would
pay for per-reading signatures, encryption and privacy-preserving encodings is
already consumed by the raw payload itself. Any scheme that adds cryptographic
protection without first reducing transmission cost will exhaust the node's
energy budget or the gateway's channel capacity.

How, then, can identity, integrity and privacy be protected simultaneously in an agricultural data
processing pipeline, end to end from constrained sensing node to distributed
ledger, within the energy, memory and channel-capacity budgets that field
deployment actually permits?

\section{Rationale of the Study}
\label{sec:rationale}

The budget for protection has to come from somewhere. If
\gls{crt} residue decomposition reduces per-packet airtime substantially, it
frees channel capacity and node energy that can be reinvested in protection
rather than merely banked as efficiency. The same decomposition has a second,
less obvious property. An individual residue discloses little about the original
value, so residues distributed across channels or custodians provide a natural
substrate for confidentiality and for threshold reconstruction. A mechanism
introduced for efficiency can thus be made to serve privacy, which is what makes
the combination worth studying rather than treating transmission and security as
separate layers.

The two mechanisms also compose cleanly. Residue
decomposition and ledger verification operate on the same object at different
points in its life, the first on the reading in flight and the second on the
record at rest, and the quantity each is sensitive to is different: the radio
cares about occupancy, the ledger about the size and count of committed
transactions. A scheme that reduces occupancy without increasing transaction
count therefore improves one without taxing the other, which is not generally
true of protection mechanisms and is worth demonstrating rather than assuming.

There is also an evidential gap. Claims about edge blockchain feasibility are
typically supported by simulation or short laboratory runs. A
sustained field deployment on real hardware, through real weather and real
connectivity failures, produces a different and more useful class of evidence,
including the failure modes that simulation does not generate.

The applied case is direct. Smallholder and arid-region agriculture is where
both the water-efficiency gains and the certification-fraud exposure are
largest, and where infrastructure assumptions of existing systems are least
likely to hold.

\section{Objectives of the Study}
\label{sec:objectives}

\subsection{Main Objective}
\label{subsec:main_objective}

\textbf{To design, implement and evaluate a distributed blockchain-based
framework that protects identity, integrity and privacy in agricultural data
processing, and to validate it in a sustained field deployment on constrained
edge hardware.}

\subsection{Specific Objectives}
\label{subsec:specific_objectives}

\begin{enumerate}
  \item \textbf{Analyse} the identity, integrity and privacy requirements of
        agricultural data pipelines, decomposing the threat surface from
        sensing node through gateway to ledger into an attacker model with
        stated capabilities and assumptions.

  \item \textbf{Design} an integrated protection scheme combining \gls{crt}
        multi-channel residue transmission, which lowers per-packet airtime
        while preserving exact reconstruction, with an identity layer that
        binds each reading to an authenticated originating node through
        \gls{hlf} \gls{msp} credentials and \gls{hrbac} policy enforced in
        chaincode.

  \item \textbf{Implement} the framework on a heterogeneous edge cluster of
        \gls{esp32} sensing nodes and \gls{arm64} \gls{rpi} gateways running
        \gls{hlf} under Raft ordering, deployed under real agricultural
        operating conditions.

  \item \textbf{Evaluate} the transmission, ledger and access-control layers
        against compression and centralised baselines on airtime, energy,
        throughput, latency and memory footprint, and assess the scalability
        and fault-tolerance envelope through queuing-theoretic analysis
        validated against measured load.
\end{enumerate}

\section{Research Questions}
\label{sec:research_questions}

\begin{enumerate}[label=\textbf{RQ\arabic*.}]
  \item To what extent can \gls{crt} residue decomposition reduce per-packet
        airtime and increase gateway packet capacity without loss of
        reconstruction fidelity?
  \item Can a permissioned distributed ledger sustain the throughput and latency
        required by real-time irrigation control when hosted entirely on
        \gls{arm64} edge gateways?
  \item What identity binding is required for a ledger record to attest to the
        originating sensing node rather than merely to the submitting gateway,
        and what does that binding cost?
  \item What confidentiality is obtained by distributing residues across
        channels and custodians, and under what adversary model does it hold?
  \item How does the integrated framework behave as offered load approaches
        gateway and ledger saturation, and where are its practical limits?
\end{enumerate}

\section{Research Hypotheses}
\label{sec:hypotheses}

\begin{description}
  \item[\textnormal{\textbf{H1 (Transmission efficiency).}}] \gls{crt} residue
        transmission reduces per-packet airtime by at least 40\% relative to raw
        transmission while maintaining at least 99\% reconstruction accuracy.

  \item[\textnormal{\textbf{H2 (Edge ledger performance).}}] \gls{hlf} hosted on
        four \gls{rpi} 4B gateways sustains a mean transaction latency below
        \SI{2}{\second} at an average throughput of at least \SI{50}{\tps}.

  \item[\textnormal{\textbf{H3 (Access-control overhead).}}] Enforcing
        \gls{rbac} in chaincode preserves at least \SI{50}{\tps} and keeps mean
        smart-contract execution below \SI{500}{\milli\second}.

  \item[\textnormal{\textbf{H4 (Field robustness).}}] Sustained field operation
        maintains at least 99\% end-to-end reliability, with measured airtime
        within \SI{5}{\milli\second} of the analytical model.

  \item[\textnormal{\textbf{H5 (Privacy preservation).}}] An adversary observing
        fewer than $k$ of the $k$ residue channels cannot reconstruct the
        original reading, and the residual information leakage remains below a
        stated bound.
\end{description}



\section{Scope and Delimitations}
\label{sec:scope}

\textbf{In scope.} Sub-GHz \gls{lora} sensing networks of the order of fifty
nodes; permissioned ledgers of the \gls{hlf} family under Raft ordering;
gateway hardware in the \gls{rpi} 4B class; agronomic telemetry, specifically
soil moisture, temperature, humidity and irrigation actuation; and the
identity, integrity and privacy properties of that pipeline.

\textbf{Out of scope.} Public permissionless blockchains and their consensus
economics; cellular and satellite backhaul; image-based and remote-sensing
agronomy; agronomic model accuracy as such, which is treated as given rather
than as an object of study; and formal cryptographic proofs of the underlying
primitives, which are cited rather than re-derived.

\textbf{Delimitations.} Field validation is drawn from a single deployment site
and a single cropping season, which bounds the claims that can be made about
generalisation across soil types and climates. Node population is bounded by
what could be instrumented. Behaviour beyond that population is established
analytically and by synthetic load rather than by direct measurement.

\section{Significance of the Study}
\label{sec:significance}

For \textbf{engineering practice}, the work supplies measured rather than
simulated evidence about what a permissioned ledger costs on edge hardware, and
a transmission scheme that reduces that cost enough to make protection
affordable.

For \textbf{the research literature}, it connects two bodies of work that have
developed separately, namely residue arithmetic for communication efficiency and
distributed ledgers for data provenance, and shows that the first can subsidise
the second.

\textbf{Methodologically}, the work contributes a characterisation obtained under
field conditions rather than in a laboratory. Failures that a testbed does not
generate, mesh partitions, thermal derating, intermittent backhaul and the
recovery behaviour that follows each, are exactly the ones that determine
whether such a system is deployable, and they only appear when the system is
left running.

For \textbf{agricultural stakeholders}, it offers a path to verifiable
certification and insurance records that does not require farmers to surrender
commercially sensitive detail to every consortium member.

\section{Definition of Key Terms}
\label{sec:definitions}

\begin{description}
  \item[Airtime] The time a radio occupies a channel to transmit one packet. It
        governs both node energy consumption and gateway capacity.

  \item[Residue] The remainder of a reading modulo one of a set of pairwise
        coprime moduli. A complete residue set uniquely determines the original
        value below the product of the moduli.

  \item[Integrity] The property that a record has not been altered since it was
        committed, and that alteration would be detectable.

  \item[Identity] The property that a record can be attributed to a specific
        authenticated principal, here a specific sensing node rather than
        merely a gateway.

  \item[Privacy] The property that a record's content is disclosed only to
        principals authorised to see it, notwithstanding replication of the
        ledger.

  \item[Endorsement] In \gls{hlf}, the process by which designated peers
        simulate a transaction and sign the result before ordering.
\end{description}

\section{Organisation of the Thesis}
\label{sec:organisation}

\textbf{Chapter One} states the background, problem, rationale, objectives,
questions and hypotheses. \textbf{Chapter Two} reviews the literature on
precision agriculture and \gls{iot} sensing, transmission optimisation,
distributed ledgers in agriculture, edge-resident consensus and
privacy-preserving data processing, and identifies the gap this work addresses.
\textbf{Chapter Three} presents the materials and methods: system architecture,
\gls{crt} formulation, identity and access-control design, hardware
implementation, deployment protocol and evaluation methodology.
\textbf{Chapter Four} reports results and applications across the transmission,
ledger, access-control and agricultural-outcome dimensions. \textbf{Chapter
Five} discusses the findings against the hypotheses, states limitations and
threats to validity, and sets out conclusions, recommendations and perspectives
for further research.
\clearpage{}
\clearpage{}

\chapter{Review of the Literature}
\label{ch:literature}

\section{Introduction}
\label{sec:lit_intro}

This chapter surveys the five bodies of work on which the thesis draws:
precision agriculture and \gls{iot} sensing; low-power wide-area transmission
and the airtime constraint that governs it; residue arithmetic and the
\gls{crt}; distributed ledgers in agriculture and their migration to the
network edge; and identity, access control and privacy for replicated records.
Each is treated on its own terms, then criticised, and the chapter closes by
naming what the five leave undone.

\section{Precision Agriculture and IoT Sensing}
\label{sec:lit_precag}

\subsection{From Wireless Sensor Networks to Managed Irrigation}

Precision agriculture began as site-specific crop management, matching inputs
to measured within-field variability rather than to field
averages~\cite{pierce1999site, zhang2002precision}. Early instrumentation used
wireless sensor networks for soil and microclimate
monitoring~\cite{wang2006wireless, bogena2010soil}, and the subsequent decade
moved the field from measurement to actuation, with \gls{lorawan} irrigation
control among the clearest applications~\cite{evans2008precision}. Reviews of
the area agree on the technical progress and on a persistent
weakness~\cite{gebbers2010precision, aydin2021comprehensive}: integration
between the sensing layer and the decision-support systems that consume it
remains fragmented, so measurement quality is rarely the limiting factor on
outcomes.

\subsection{What the Sensing Literature Does Not Settle}

Two questions are left open. First, almost all of this work treats the reading
as trustworthy by construction: the sensor reports, the system believes it, and
no mechanism attests that the value committed to storage is the value measured.
Second, the economics that make a deployment viable on a smallholding, which
force the node bill of materials to a few tens of dollars, are seldom carried
through to the security design. A protection mechanism that assumes a secure
element or a mains supply has assumed away the constraint that defines the
problem.

\section{Low-Power Wide-Area Transmission and the Airtime Constraint}
\label{sec:lit_lora}

\subsection{The Airtime Model}

\gls{lora} is the dominant physical layer for agricultural telemetry, with
demonstrated rural ranges of several kilometres and, in favourable terrain, up
to \SI{15}{\kilo\metre}~\cite{adelantado2017understanding, li2020performance,
raj2017comprehensive}. Its cost is airtime. For a spreading factor $SF$,
bandwidth $BW$ and coding rate $CR$, the symbol duration and the time on air of
a packet carrying $PL$ payload bytes are

\begin{equation}
  T_{\text{sym}} = \frac{2^{SF}}{BW},
  \label{eq:symbol_time}
\end{equation}

\begin{equation}
  T_{\text{air}} = T_{\text{preamble}}
    + T_{\text{sym}}\!\left(8 + \max\!\left(
      \left\lceil \frac{8PL - 4SF + 28 + 16}{4(SF - 2DE)} \right\rceil
      (CR + 4),\; 0 \right)\right),
  \label{eq:airtime}
\end{equation}

where $DE$ is the low-data-rate optimisation flag. The term that matters is the
ceiling: airtime grows in discrete steps with payload length, and every
additional symbol is charged at $T_{\text{sym}}$, which doubles with each
increment of $SF$~\cite{adelantado2017understanding, bor2016lorawan}.

Airtime is not merely an energy cost. Sub-GHz bands impose a duty-cycle limit,
so a node that has transmitted for $T_{\text{air}}$ must remain silent for a
proportional interval, which caps reporting frequency independently of battery
capacity. Under the unslotted access that \gls{lorawan} class A uses, the
offered load per channel over a reporting window $T_w$ for $N$ nodes is

\begin{equation}
  G = \frac{N \cdot T_{\text{air}}}{T_w},
  \label{eq:offered_load}
\end{equation}

and the pure-ALOHA collision probability follows as

\begin{equation}
  P_{\text{coll}} = 1 - e^{-2G} \approx 2G \quad \text{for small } G,
  \label{eq:collision}
\end{equation}

which is linear in $T_{\text{air}}$ in the low-load regime and degrades sharply
as $G$ approaches unity~\cite{bor2016lorawan, adelantado2017understanding}.
Reducing per-packet airtime therefore buys capacity twice over, once through
the duty cycle and once through the collision term.

\subsection{Compression and Sampling Responses}

The conventional response is to shorten the payload. Lossless coders adapted to
sensor streams, principally \gls{lzw} and Huffman variants, achieve useful
ratios on agronomic data~\cite{li2018lossless, marcelloni2009energy}, but their
dictionary and buffer requirements are measured in kilobytes, which is a
significant fraction of the free \gls{ram} on a microcontroller already running
a radio stack and a sensor driver~\cite{chen2025compression}. Adaptive sampling
transmits only when a reading departs from a predicted
trajectory~\cite{kjellsson2008adaptive}, at the risk of suppressing exactly the
transient that mattered. Aggregation and protocol-level
optimisation~\cite{heinzelman2000energy, sanchez2021energy} reduce traffic at
the cost of resolution or latency.

\subsection{The Common Limitation}

Every one of these approaches reduces the number of bytes sent. None reduces
the occupancy of an individual packet while preserving the reading exactly, and
occupancy is what Equations~\ref{eq:offered_load} and~\ref{eq:collision} are
sensitive to. This distinction, between total volume and per-packet occupancy,
is where the present work departs from the compression literature.

\section{Residue Arithmetic and the Chinese Remainder Theorem}
\label{sec:lit_crt}

\subsection{Statement and Reconstruction}

Let $m_1, m_2, \ldots, m_k$ be pairwise coprime moduli and
$M = \prod_{i=1}^{k} m_i$. For any integer $x$ with $0 \le x < M$, the residue
vector

\begin{equation}
  \mathbf{r} = (r_1, r_2, \ldots, r_k), \qquad r_i = x \bmod m_i,
  \label{eq:residues}
\end{equation}

determines $x$ uniquely, and $x$ is recovered by

\begin{equation}
  x = \left( \sum_{i=1}^{k} r_i M_i N_i \right) \bmod M,
  \qquad M_i = \frac{M}{m_i}, \quad N_i = M_i^{-1} \bmod m_i .
  \label{eq:crt_reconstruction}
\end{equation}

Reconstruction is exact because $M_i N_i \equiv 1 \pmod{m_i}$ and
$M_i \equiv 0 \pmod{m_j}$ for $j \neq i$~\cite{hardy2008introduction,
ding1996remainder}. Evaluated by Garner's algorithm the cost is $O(k)$ modular
operations with no division in the inner loop~\cite{garner1959residue}, which is
what makes the scheme viable on a gateway rather than only on a workstation.
Figure~\ref{fig:crt-concept} shows the decomposition and its inverse.

\begin{figure}[htbp]
  \centering
  \begin{tikzpicture}[
      b/.style={draw, rounded corners=2pt, align=center, font=\small,
                minimum height=8mm, inner sep=3pt},
      r/.style={b, fill=black!5, text width=17mm},
      a/.style={-{Stealth[length=2mm]}, thick},
      lbl/.style={font=\footnotesize\itshape, align=center}
    ]
    \node[b, text width=22mm] (x)  at (0,0) {reading\\$x < M$};
    \node[r] (r1) at (3.6, 1.4) {$r_1 = x \bmod m_1$};
    \node[r] (r2) at (3.6, 0)   {$r_2 = x \bmod m_2$};
    \node[r] (r3) at (3.6,-1.4) {$r_3 = x \bmod m_3$};
    \node[b, text width=24mm] (rec) at (7.8,0) {Garner\\reconstruction\\$O(k)$};
    \node[b, text width=22mm] (y)  at (11.4,0) {$x$ recovered\\exactly};

    \foreach \n in {r1,r2,r3} { \draw[a] (x.east) -- (\n.west);
                                \draw[a] (\n.east) -- (rec.west); }
    \draw[a] (rec) -- (y);
    \node[lbl] at (3.6,-2.4) {three short packets,\\demodulated concurrently};
  \end{tikzpicture}
  \caption{Residue decomposition and reconstruction. One long packet is
           replaced by $k$ short ones that a multi-channel concentrator can
           demodulate simultaneously, and the original value is recovered
           exactly in $O(k)$ operations.}
  \label{fig:crt-concept}
\end{figure}

\subsection{Where Residue Systems Have Been Applied}

Residue number systems built on Equation~\ref{eq:crt_reconstruction} are long
established in computer arithmetic, where carry-free addition and
multiplication across residue channels accelerate arithmetic and support
fault-tolerant computation~\cite{parhami2010computer, shoup2009computational}.
The theorem is equally standard in cryptography~\cite{ding1996remainder,
garner1959residue} and in error correction~\cite{goldreich1999china}. Within
sensor networks, \textcite{wang2018novel} applied residue representations to
reduce stored data volume.

\subsection{The Gap in the Residue Literature}

The sensor-network application targeted storage, not transmission, and
therefore did not exploit the property that matters here: that $k$ residues are
$k$ independent short packets which a multi-channel concentrator can receive at
the same instant. The computer-arithmetic literature exploits parallelism, but
across arithmetic units rather than across radio channels, where the currency
is occupancy rather than processor cycles. No prior system, to the author's
knowledge, uses residue decomposition to drive parallel demodulation on a
low-power radio.



\section{Distributed Ledgers in Agriculture}
\label{sec:lit_blockchain}

\subsection{Supply-Chain Transparency}

The earliest and best-developed application is provenance along the food
supply chain. \textcite{tian2016agri} combined RFID with a blockchain to trace
agri-food products end to end, and subsequent surveys map the
field~\cite{kamilaris2019rise, galvez2018future} and its food-safety
consequences~\cite{oh2025foodsafety}. The common architecture places the ledger
between organisations that have commercial reason to distrust one another,
which is precisely the setting in which replicated tamper-evident storage earns
its cost.

\subsection{Farm Data Management and Smart-Contract Automation}

A second line applies the ledger inside the farm, to data ownership and access
rather than to product movement~\cite{lin2018blockchain, ferrag2020security,
haque2024scalable}. A third automates agreements, using smart contracts to
trigger payments, insurance settlement or irrigation
actions~\cite{rama2021systematic, patil2022blockchain}. Both assume a ledger
that is reachable and responsive.

\subsection{The Recurring Weakness}

Across all three lines, implementations that operate under the constraints of
real-time control on resource-constrained hardware remain scarce. Where
performance is reported it is typically obtained on server-class hosts, and the
question of whether endorsement and ordering complete inside an irrigation
control loop is not usually asked.

\section{Edge-Resident Consensus and Constrained Hardware}
\label{sec:lit_edge}

\subsection{Edge Computing in Agricultural Systems}

Moving computation to the field reduces round-trip latency and reduces
dependence on backhaul that rural sites often lack~\cite{sharma2020edge}.
\textcite{liu2019edge} reported improved response times for irrigation control
under edge processing, and \textcite{wang2021edge} showed that keeping data
local improves privacy by limiting what leaves the site. \gls{arm64}
single-board clusters are now routinely used for distributed sensing and
inference~\cite{sharma2020edge, wang2021edge}.

\subsection{Networking the Edge}

Agricultural deployments overwhelmingly use star topologies, in which a single
gateway failure removes a zone. Mesh routing has been studied mostly in urban
and industrial settings~\cite{akylidiz2005wireless}; BATMAN-adv has shown
promise for rural connectivity~\cite{schmidt2010b.a.t.m.a.n.} without being
taken up in precision agriculture. Hybrid multi-technology architectures have
been proposed~\cite{chettri2020comprehensive} but tend to add management
complexity without integrating cleanly.

\subsection{Capacity Limits}

Where a gateway serves $C$ channels over a window $W$ and each node consumes
$T_{\text{eff}}$ of effective airtime per reporting cycle, the node ceiling is

\begin{equation}
  N_{\max} = \frac{C \cdot W}{T_{\text{eff}}},
  \label{eq:node_ceiling}
\end{equation}

and per-channel packet throughput, for a concentrator that demodulates one
packet per channel concurrently, is

\begin{equation}
  \lambda = \frac{1}{T_{\text{air}}}.
  \label{eq:packet_rate}
\end{equation}

Equations~\ref{eq:node_ceiling} and~\ref{eq:packet_rate} pull in different
directions once a reading is split across channels: splitting raises the total
airtime a node consumes, and so lowers $N_{\max}$, while lowering
$T_{\text{air}}$ per packet, and so raising $\lambda$. Which of the two binds
depends on whether the deployment is limited by node count or by packet
capacity, a distinction the existing literature does not draw.

\subsection{What Is Missing at the Edge}

Edge implementations rarely integrate with verifiable storage, so the trust gap
that motivates the ledger reappears the moment computation is moved off the
ledger. Communication parallelism, meaning simultaneous use of multiple
wireless channels by one logical reading, is essentially unexplored.

\section{Identity, Access Control and Privacy}
\label{sec:lit_privacy}

\subsection{Identity and Access Control}

Permissioned ledgers supply identity through a \gls{msp} backed by a
\gls{pki}, which allows a record to be attributed to an enrolled principal
rather than merely to an address~\cite{tar2023fabric, jayaram2023hyperledger}.
\gls{rbac} and its hierarchical variants are the usual policy layer, enforced
in chaincode so that the authorisation decision is itself
auditable~\cite{ferrag2020security}.

\subsection{Privacy Under Replication}

Replication is what makes a ledger tamper-evident and is also what makes it
leak. Where a consortium includes buyers, insurers and cooperatives, every
member sees every replicated record, and agronomic traces disclose plot
productivity and irrigation strategy. Keeping data local reduces the exposure
but reintroduces the trusted party the ledger was meant to
remove~\cite{wang2021edge}.

\subsection{The Unresolved Tension}

The literature offers identity and integrity together, and privacy separately,
usually by withholding data from the ledger. What it does not offer is a
mechanism in which the same operation that makes transmission affordable also
constrains disclosure. That possibility, latent in the residue representation
of Equation~\ref{eq:residues}, is the opening this thesis works in.

\section{Research Gaps}
\label{sec:lit_gaps}

Table~\ref{tab:gaps} states the gaps that follow from the five preceding
sections and the objective that addresses each.

\begin{table}[htbp]
  \centering
  \caption{Gaps identified in the literature and the objective addressing each.}
  \label{tab:gaps}
  \begin{singlespace}
  \small
  \begin{tabular}{@{}p{3.1cm}p{6.6cm}p{2.6cm}@{}}
    \toprule
    \textbf{Area} & \textbf{Gap} & \textbf{Objective} \\
    \midrule
    Transmission cost
      & Compression reduces total volume, not per-packet occupancy, and costs
        memory a constrained node does not have
      & O2 \\
    Residue arithmetic
      & Applied to storage and to arithmetic units, never to parallel
        demodulation on a low-power radio
      & O2 \\
    Edge ledgers
      & Performance characterised on server hardware; behaviour on \gls{arm64}
        gateways under partition unmeasured
      & O3, O4 \\
    Identity
      & Authentication terminates at the gateway, so the record attests to the
        submitter rather than the sensor
      & O1, O2 \\
    Privacy
      & Confidentiality obtained by withholding data, which reinstates a
        trusted intermediary
      & O1, O2 \\
    Capacity
      & No validated model distinguishing node-count from packet-capacity
        limits for multi-channel residue traffic
      & O4 \\
    \bottomrule
  \end{tabular}
  \end{singlespace}
\end{table}

\section{Partial Conclusion}
\label{sec:lit_conclusion}

The sensing literature establishes what to measure and treats the measurement
as trustworthy. The transmission literature establishes that airtime binds and
attacks it by shortening payloads, which does not address per-packet occupancy.
Residue arithmetic supplies a decomposition whose parallelism has been
exploited in silicon but not on the air. The ledger literature supplies
tamper evidence and identity, characterised on hardware that agricultural
deployments do not have. The privacy literature obtains confidentiality by
withholding, which forfeits verifiability.

The four therefore leave a single composite question, which the remainder of
this thesis addresses: whether a decomposition adopted to reduce per-packet
occupancy can simultaneously carry node-level identity and bound disclosure,
and whether the resulting system holds up on the hardware and in the conditions
that field deployment actually imposes.
\clearpage{}


\ubfrontheading{References}
\begin{singlespace}
\printbibliography[heading=none]
\end{singlespace}

\appendix
\titleformat{\chapter}[display]
  {\normalfont\bfseries\centering\singlespacing\fontsize{16}{18.4}\selectfont}
  {\MakeUppercase{Appendix \thechapter}}
  {0pt}
  {\MakeUppercase}
\clearpage{}

\chapter{Supplementary Material}
\label{app:supplementary}

\clearpage{}


\end{document}
